%------------------------------------------------------------------------------%
\begin{question}
  Dada a função
  \[
    f(x, y) = (x-2)^2 + (y-1)^2
  \]

  a) Determine e esboce a curva de nível \(f(x, y) = 4\)

  b) Calcule e esboce o gradiente de $f$ nos pontos
  \((0,  1)\),
  \((2, -1)\),
  \((2,  3)\) e
  \((4,  1)\)
\end{question}

%------------------------------------------------------------------------------%
\begin{resolution}{Solução}
  \onslide<+->{a) Determinando a curva de nível}
  \begin{align*}
    \onslide<+->{f(x, y)           & = 4}   \\
    \onslide<+->{(x-2)^2 + (y-1)^2 & = 2^2}
  \end{align*}
  \onslide<+->{Equação da circunferência centrada em \((2, 1)\) e raio 2}
\end{resolution}

%------------------------------------------------------------------------------%
\begin{resolution}{Solução}
  b) Calculando as derivadas parciais de $f$

  \pause
  \[
    \D{f}{x}
    \pause
    \;=\; \D{}{x}\left[(x-2)^2 + (y-1)^2\right]
    \pause
    \;=\; 2(x-2)\D{}{x}(x-2)
    \pause
    \;=\; 2(x-2)
  \]

  \pause
  \[
    \D{f}{y}
    \pause
    \;=\; \D{}{y}\left[(x-2)^2 + (y-1)^2\right]
    \pause
    \;=\; 2(y-1)\D{}{y}(y-1)
    \pause
    \;=\; 2(y-1)
  \]

  \pause
  Gradiente de $f$
  \pause
  \[
    \nabla f(x, y)
    \pause
    \;=\; \nabla\left[(x-2)^2 + (y-1)^2\right]
    \pause
    \;=\; \begin{pmatrix}
      2(x-2) \\ 2(y-1)
    \end{pmatrix}
  \]
\end{resolution}

%------------------------------------------------------------------------------%
\begin{resolution}{Solução}
  Calculando o gradiente nos pontos
  \((0,  1)\) e
  \((2, -1)\)

  \pause
  \[
    \nabla f(0, 1)
    \pause
    \;=\; \begin{pmatrix}
      2(x-2) \\ 2(y-1)
    \end{pmatrix} \evalat{(0, 1)}{}
    \pause
    \;=\; \begin{pmatrix}
    2(0-2) \\ 2(1-1)
    \end{pmatrix}
    \pause
    \;=\; \begin{pmatrix}
    -4 \\ 0
    \end{pmatrix}
  \]

  \pause
  \[
    \nabla f(2, -1)
    \pause
    \;=\; \begin{pmatrix}
      2(x-2) \\ 2(y-1)
    \end{pmatrix} \evalat{(2, -1)}{}
    \pause
    \;=\; \begin{pmatrix}
      2(2-2) \\ 2(-1-1)
    \end{pmatrix}
    \pause
    \;=\; \begin{pmatrix}
      0 \\ -4
    \end{pmatrix}
  \]
\end{resolution}

%------------------------------------------------------------------------------%
\begin{resolution}{Solução}
  Calculando o gradiente nos pontos
  \((2,  3)\) e
  \((4,  1)\)

  \pause
  \[
    \nabla f(2, 3)
    \pause
    \;=\; \begin{pmatrix}
      2(x-2) \\ 2(y-1)
    \end{pmatrix} \evalat{(2, 3)}{}
    \pause
    \;=\; \begin{pmatrix}
      2(2-2) \\ 2(2-1)
    \end{pmatrix}
    \pause
    \;=\; \begin{pmatrix}
      0 \\ 4
    \end{pmatrix}
  \]

  \pause
  \[
    \nabla f(4, 1)
    \pause
    \;=\; \begin{pmatrix}
      2(x-2) \\ 2(y-1)
    \end{pmatrix} \evalat{(4, 1)}{}
    \pause
    \;=\; \begin{pmatrix}
      2(4-2) \\ 2(1-1)
    \end{pmatrix}
    \pause
    \;=\; \begin{pmatrix}
      4 \\ 0
    \end{pmatrix}
  \]
\end{resolution}

%------------------------------------------------------------------------------%